\documentclass[12pt]{article}
\usepackage{amsmath,setspace,graphicx,amssymb,amsthm,float}
\usepackage{hyperref}
\hypersetup{
    colorlinks,
    allcolors=black
}
\title{\textbf{MATH2102 Real Analysis Lecture Notes}}
\author{Thomas Piercy}
\date{\today{}}

\newcommand{\df}[2]{\frac{\text{d}#1}{\text{d}#2}}
\newcommand{\seq}[2]{(#1_#2)_{#2\in\mathbb{N}}}
\newcommand{\pf}[2]{\frac{\partial #1}{\partial #2}}
\begin{document}
\maketitle
\hrulefill
\tableofcontents
\newpage
\setstretch{1.5}
\section{Sequences, series, and limits}
\subsection{Sequence Convergence Definition}
A sequence \(\seq{a}{n}\) is said to converge to a limit \(L\), if \(\forall\,\varepsilon>0\) there exists an \(N\in\mathbb{N}\) such that \(\forall\,n\in\mathbb{N}\) with \(n\geq N\) there holds: 
\[|a_n-L|<\varepsilon\]
We write \(a_n\to L\) as \(n\to \infty\) or \(\lim_{n\to\infty}a_n=L\)
\subsubsection*{Example}
\[a_n=\frac{1}{\sqrt{n}}\]
We want to show \(a_n\to 0\) as \(n\to\infty\). Let \(\varepsilon>0\). Define \(N=\lfloor \frac{1}{\varepsilon^2}\rfloor+1\). Then for \(n\geq N\):
\[|a_n-0|=\frac{1}{\sqrt{n}}\leq\frac{1}{\sqrt{N}}<\varepsilon\]
\[\blacksquare\]
\subsubsection{Convergence to \(\infty\)}
A sequence \(\seq{a}{n}\) converges to \(\infty\) if \(\forall K\in\mathbb{R}\) there exists \(N\in\mathbb{N}\) such that \(\forall n\geq N\) there holds \(a_n>K\). We write \(a_n\to\infty\) as \(n\to\infty\).
\subsection{Convergence of Reciprocals}
Let \(\seq{a}{n}\) be a sequence. If \(a_n>0,\ \forall n\in\mathbb{N}\), then \(a_n\to 0\iff \frac{1}{a_n}\to\infty\) as \(n\to\infty\).
\subsection{Algebra of Limits}
Similar to Core Mathematics in Year 1, we will make great use of algebra of limits in problem solving. A quick recap is that for sequences \(\seq{a}{n}\) and \(\seq{b}{n}\) with \(a_n\to a\) and \(b_n\to b\) as \(n\to \infty\) then we have the results:
\begin{itemize}
    \item \(a_n+b_n\to a+b\)
    \item \((a_n)(b_n)\to ab\)
    \item \(\frac{1}{a_n}\to\frac{1}{a}\) subject to some conditions: namely \(a_n\ne 0\) for all \(n\geq N\in\mathbb{N}\).
\end{itemize}
\subsection{Limits of Inequalities, Squeeze Theorem}
If \(\seq{a}{n}\) and \(\seq{b}{n}\) are sequences with \(a_n\leq b_n\) for all \(n\) and \(a_n\to a\), \(b_n\to b\) as \(n\to\infty\) then \(a\leq b\). If \(\seq{a}{n},\,\seq{b}{n}\) and \(\seq{c}{n}\) are sequences with \(a_n\leq b_n\leq c_n\) for all \(n\) such that \(a_n\to L\) and \(c_n\to L\) as \(n\to\infty\), then we also have \(b_n\to L\) as \(n\to \infty\).
\subsubsection*{Example}
Let \(\seq{a}{n}=\frac{2^n}{n!}\). Then we have:
\[0\leq \frac{2^n}{n!}\leq \frac{4}{n}\]
However we know that \(\frac{4}{n}\to 0\) as \(n\to\infty\), then by extension we have \(a_n\to 0\) as \(n\to \infty\) by the squeeze theorem.
\subsection{Ratio Test for Sequences}
Let \(\seq{a}{n}\) be a sequence with \(a_n\geq 0\) for all \(n\). Assume the sequence of ratios:
\[r_n=\frac{a_{n+1}}{a_n}\]
converges to \(q\in[0,\infty]\). We have the following result:
\[a_n\to\begin{cases}
    0 & 0\leq q<1 \\
    \infty & q>1 \\
    \text{indeterminate} & q=1
\end{cases}\]
\subsubsection{Proof for \(q>1\) case}
Let \(\seq{a}{n}\) be a sequence such that:
\[\frac{a_{n+1}}{a_n}\to q>1\]
Define a new value:
\[s=\frac{1+q}{2}\]
such that \(q>s>1\). Set:
\[\varepsilon=\frac{q-1}{2}>0\]
Then for \(n>N\) we have:
\[\left|\frac{a_{n+1}}{a_n}-q\right|<\varepsilon\]
\[\implies \frac{a_{n+1}}{a_n}>s>1\]
\[\implies \frac{a_n}{a_N}>s^{n-N}\]
\[\implies a_n>(a_Ns^{-N})s^n\]
So from squeeze theorem, we have \(a_n\to \infty\) as \(n\to \infty\).
\subsection{Standard Sequences}
These sequences should be remembered and can be quoted without proof in an exam. Let \(K\in\mathbb{Z}\) then we have the following results:
\[n^k\to\begin{cases}
    0 & k<0 \\
    1 & k=0 \\
    \infty & k>0
\end{cases}\]
Let \(q>0\):
\[q^n\to\begin{cases}
    0 & q<1 \\
    1 & q=1 \\
    \infty & q>1
\end{cases}\]
Let \(q\leq 0\):
\[q^n\to\begin{cases}
    0 & q>-1 \\
    \text{diverges} & \text{otherwise}
\end{cases}\]
Let \(0<q<1\) and \(K>0\):
\[n^kq^n\to 0,\quad n\to\infty\]
\subsubsection*{Example}
Let \(\seq{a}{n}\) be a sequence with:
\[a_n=\frac{2^n}{n!}\]
Then we have:
\[\frac{a_{n+1}}{a_n}=\left(\frac{2^{n+1}}{(n+1)!}\right)\left(\frac{n!}{2^n}\right)=\frac{2}{(n+1)}\to 0\text{ as }n\to\infty\]
\[\therefore \frac{2^n}{n!}\to 0\text{ as }n\to\infty\]
\subsection{Monotonicity}
A sequence \(\seq{a}{n}\) is called monotone increasing if for all \(n\in\mathbb{N}\), we have \(a_n\leq a_{n+1}\). It is called monotone decreasing if for all \(n\in\mathbb{N}\), we have \(a_n\geq a_{n+1}\). Constant sequences are monotone. Sequences like \((-1)^n\) are neither increasing nor decreasing. If \(a_n > 0\) for all \(n\) then \(\seq{a}{n}\) is monotone increasing if and only if \(a_{n+1}/a_n\geq 1\) for all \(n\).
\subsection{Supremum and Infimum}
Let \(A\subseteq \mathbb{R}\) be a non-empty set. If \(A\) is bounded above it has a least upper bound \(\sup A\in \mathbb{R}\). We have \(s=\sup A\) if:
\begin{align*}
    &\forall x\in A,\ x\leq s \\
    &\forall \varepsilon>0,\, \exists y\in A:\ y>s-\varepsilon
\end{align*}
If \(A\) is unbounded above, we write \(\sup A=+\infty\). If \(A\) is bounded below then it has a greatest lower bound \(\inf A\in\mathbb{R}\). We have \(l=\inf A\) if:
\begin{align*}
    &\forall x\in A,\ x\geq l \\
    &\forall \varepsilon>0,\, \exists y\in A:\ y<l+\varepsilon
\end{align*}
If \(A\) is unbounded below, we write \(\inf A=-\infty\). We set:
\begin{align*}
    &\sup \emptyset =-\infty \\
    &\inf \emptyset = +\infty
\end{align*}
For \(A\subseteq \mathbb{R}\) and \(A\ne\emptyset\), we have \(\inf A\leq \sup A\). If \(A\subseteq B\subseteq \mathbb{R}\) then \(\sup A\leq \sup B\) and \(\inf A \geq \inf B\). A monotone increasing sequence converges to it's supremum and a monotone decreasing sequence converges to it's infimum. A bounded monotone sequence always converges.
\end{document}